%%%%%%%%%%%%%%%%%%%%%%%%%%%%%%%%%
%
%       EXERCÍCIO
%
%%%%%%%%%%%%%%%%%%%%%%%%%%%%%%%%%

\ifdefstring{\atividade}{prova}{%
    \ifdefstring{\modo}{objetivo}{%
        \renewcommand{\valorquestao}{\ValorQObj\ ponto}
    }{%
        \renewcommand{\valorquestao}{\ValorQDisc\ pontos}
    }%
}%

\begin{exercicioBanco}[\valorquestao]
Considere as afirmações a seguir e responda se saõ verdadeiras ou falsas.
\begin{enumerate}[(i)]
    \item Seja \(k \in \bbR\). O conjunto solução da inequação \(|x| < k\) é \(S = ]-\infty,-k[\;\cup\;]k,\infty[\).
    %
    \item Para todo \(x, y \in \bbR\) vale que \(|x + y| = |x| + |y|\).
    %
    \item \(|-2x| = -2|x|\), para todo \(x \in \bbR\).
\end{enumerate}
% Define as alternativas
\newcommand{\alternativas}{%
    \begin{center}
        \begin{tabularx}{\textwidth}{XXXXX}
            (a) (F, F, V). &
            (b) (F, V, V). &
            (c) (V, F, F). &
            (d) (F, V, F). &
            (e) (F, F, F).
        \end{tabularx}
    \end{center}
}

% Define a resposta correta
\newcommand{\resposta}{E}

% Lógica condicional para exibição
\ifdefstring{\atividade}{lista}{%
    \alternativas
    \vspace{0.5em}
    
    \noindent\textbf{Resposta:} letra \textbf{\resposta}.
}{%
    \ifdefstring{\modo}{objetivo}{%
        \alternativas
    }{%
        % modo = discursiva → não mostra alternativas
    }
}
\end{exercicioBanco}

